\subsection{Crystallographic Restriction}
\begin{theorem}
    \textbf{Crystallographic restriction}:

    Als een kristal invariant is onder \(\frac{2\pi }{n}\) rotaties rond een as, dan moet \(n\) gelijk zijn aan \(1, 2, 3, 4\) of \(6\).
    
\end{theorem}


\begin{proof}
   
    Ruimtegroepen hebben een extra voorwaarde nog eens op de groepen \(SO_3\) en \(O_3\) die we al hebben gezien,
    namelijk dat de elementen \(O_g\) equivalent moeten zijn aan integer matrices in een geschikte basis.
    Waar de matrix van de basisvectoren \(E^{-1}O_gE = N_g\) is.
     

    De eigenwaarden van een integer matrix corresponderend aan een element met orde \(n\) zijn \(\pm 1, \pm e^{\frac{i2\pi}{n}}, \pm e^{-\frac{i2\pi}{n}} \).
    
    Doordat de matrix \(O_g\) equivalent is aan een integer matrix \((N_g)\), moet de trace van \(O_g\) een geheel getal zijn.
    
    De matrix \(O_g\) is:
    \[
        O_g = \begin{pmatrix}
            1 & 0 & 0 \\
            0 & \cos(\frac{2\pi}{n}) & -\sin(\frac{2\pi}{n}) \\
            0 & \sin(\frac{2\pi}{n}) & \cos(\frac{2\pi}{n})
        \end{pmatrix}
    \]
    \(Tr (O_g) = 1 + 2\cos(\frac{2\pi}{n})\).

    De voorwaarde dat de trace een geheel getal is en dat \(\cos \,x , \quad\forall x \in [1,-1] \implies Tr(O_g) \in \{-1,0,1,2,3\} \),
    wat overeenkomt met \(n = 1, 2, 3, 4\) of \(6\).
    \begin{itemize}
        \item[\(n = 1\):] \(1 + 2\cos(\frac{2\pi}{1}) = 3\)
        \item[\(n = 2\):] \(1 + 2\cos(\frac{2\pi}{2}) = -1\)
        \item[\(n = 3\):] \(1 + 2\cos(\frac{2\pi}{3}) = 0\)
        \item[\(n = 4\):] \(1 + 2\cos(\frac{2\pi}{4}) = 1\)
        \item[\(n = 6\):] \(1 + 2\cos(\frac{2\pi}{6}) = 2\)
    \end{itemize}
\end{proof}